\section{Lesson 06 --- 20/10/2026}

\subsection{Some notation}
From now on we will use $D$ to denote the training set (the subset of training tuples reaching the current node) and $|D|$ for the cardinality of $D$ (number of tuples). We consider $m$ classes $C_1,\dots,C_m$ and we call $p_i$ the empirical probability that a tuple in $D$ belongs to class $C_i$:
\[
  p_i = \frac{|\{x \in D : y(x) = C_i\}|}{|D|}.
\]
Then we can define the \g{entropy} as the expected information needed to classify a tuple in $D$. The lower the entropy, the purer $D$ is (it is $0$ when all the tuples belong to the same class).

\begin{definition}[Entropy]
  \[
    \mathit{Info}(D) = -\sum_{i=1}^{m} p_i \log_2 p_i .
  \]
\end{definition}

A candidate attribute $A$ with $v$ different outcomes $a_1,\dots,a_v$ will partition $D$ into $v$ subsets $D_j = \{x \in D : A(x) = a_j\}$, $j = 1,\dots,v$. We can then measure the \g{weighted entropy} after splitting on $A$ as the weighted average of the entropies of the child nodes:
\[
  \mathit{Info}_A(D) = \sum_{j=1}^{v} \frac{|D_j|}{|D|}\, \mathit{Info}(D_j).
\]

When choosing a partition we need to be careful not to choose partitions that lead to a single instance in each node: the entropy will be $0$ because every child is a pure partition, but the split is useless for classifying new data.

The \g{information gain} is the reduction in entropy after the split.

\begin{definition}[Information gain]
  \[
    \mathit{Gain}(A) = \mathit{Info}(D) - \mathit{Info}_A(D).
  \]
\end{definition}

We want to maximize $\mathit{Gain}(A)$, so at each node we select the attribute with the highest information gain (this is what ID3 does).

\subsubsection{Split-point search}
For a continuous attribute $A$ we cannot create one branch per value, so we look for a threshold $s$ and make a binary split ($v = 2$):
\[
  D_1(s) = \{x \in D : A(x) \le s\}, \qquad D_2(s) = \{x \in D : A(x) > s\}.
\]
We sort the observed values of $A$ in increasing order, $a_1 \le a_2 \le \dots \le a_n$. The candidate thresholds are the midpoints between adjacent values,
\[
  s_\ell = \frac{a_\ell + a_{\ell+1}}{2}, \qquad \ell = 1,\dots,n-1,
\]
so there are at most $n-1$ candidates. For each of them we compute the weighted entropy of the two children, and we keep the threshold that minimizes it (that is, the one with the highest gain):
\[
  s^* = \argmin_{s} \mathit{Info}_{A,s}(D).
\]
The node then tests $A \le s^*$ versus $A > s^*$.

\subsection{Decision tree induction: top-down greedy learning}

At each node the algorithm chooses the best local split and recursively partitions the data (see the algorithm in the previous lesson):
\begin{enumerate}
  \item all the training instances start at the root;
  \item we select a test attribute using a heuristic or statistical measure (information gain, gain ratio, Gini index, $\chi^2$, \dots);
  \item each child receives the subset $D_j$ of instances satisfying its branch condition, and we repeat recursively;
  \item we stop when all the instances have the same class, no attributes remain or no instances are left.
\end{enumerate}
A partition is pure when all its tuples belong to the same class: good splits produce children that are purer than the parent node.

\subsubsection{Gain ratio}
Information gain can favor attributes with a lot of different outcomes, such as unique identifiers, that create many tiny, apparently pure partitions. C4.5, the successor of ID3, uses the \g{gain ratio} to prevent that: it normalizes the information gain by the \g{split information}, the information generated by the split itself.

\begin{definition}[Gain ratio]
  \[
    \mathit{SplitInfo}_A(D) = -\sum_{j=1}^{v} \frac{|D_j|}{|D|} \log_2 \frac{|D_j|}{|D|},
    \qquad
    \mathit{GainRatio}(A) = \frac{\mathit{Gain}(A)}{\mathit{SplitInfo}_A(D)}.
  \]
\end{definition}

We then select the attribute with the maximum gain ratio.

\begin{remark}
  The gain ratio can prefer highly unbalanced splits: when $\mathit{SplitInfo}_A(D)$ is close to $0$ the ratio becomes very large.
\end{remark}

% ToDo: Gini(D)

\subsection{Overfitting}

If we have noise and irregularities in the training set and the classifier is too tuned on the training set, it will also be tuned on the noise: this results in \g{overfitting} and in a decrease of the generalization performance. The overfitted tree is often large, with too many branches. To avoid overfitting we need to reduce the complexity of the classifier; of course, if it is too simple we can have \g{underfitting} and not be able to classify well, so we need a compromise.

To detect overfitting we can follow the assessment procedure we already discussed in the previous lessons: compare the training accuracy with the accuracy on an independent test (or validation) set. The test set must not be used to build the tree, otherwise the estimate of the generalization performance is biased.

Some notation:
\begin{itemize}
  \item \g{re-substitution error}: the error computed on the same training set used to learn the tree. It is often too optimistic, because the tree has been fitted to these examples;
  \item \g{generalization error}: the error expected on independent, previously unseen samples. It is the quantity we really want to minimize;
  \item \g{validation} (or \g{pruning}) \g{set}: a separate dataset used only to estimate whether pruning a subtree is beneficial.
\end{itemize}

For a tree $T$ evaluated on $n$ samples:
\[
  e_{TR}(T) = \frac{\#\text{training misclassifications}}{n_{TR}},
  \qquad
  e_{TS}(T) = \frac{\#\text{test or validation misclassifications}}{n_{TS}}.
\]
The optimistic estimate takes $e_{TS}(T) \approx e_{TR}(T)$. With reduced error pruning we use a pruning set to estimate $e_{TS}(T)$ directly, while the pessimistic estimate penalizes the training error by the number of leaves.

\subsection{Pruning}

We can use pruning to reduce the tree complexity and overfitting. There are two different strategies: \g{pre-pruning} and \g{post-pruning}.
\begin{itemize}
  \item In pre-pruning we prune the tree during its construction: we stop splitting when some stopping criterion is met (for example maximum depth, minimum number of samples per node, minimum impurity decrease), so we need to fix some thresholds. The main problem is that we could stop too early and underfit.
  \item In post-pruning we prune after growing a large tree, by replacing subtrees with leaves using an estimated-error criterion. The new leaf is labelled with the majority class among the examples reaching that node. The main limitation is that it requires growing and evaluating a large tree.
\end{itemize}
Usually post-pruning is safer than pre-pruning, because potentially useful splits are not rejected too early.

\subsubsection{Reduced error pruning}
In \g{reduced error pruning} (REP) we use a separate pruning set to evaluate whether a subtree needs to be removed or not. Let $T_v$ be the subtree rooted at node $v$; on the pruning set we compute
\begin{itemize}
  \item $\mathit{misclass}(T_v)$: number of errors made by the subtree;
  \item $\mathit{misclass}(v)$: number of errors if $v$ is replaced by a leaf;
\end{itemize}
and then
\[
  \mathit{Gain\_prune}(v) = \mathit{misclass}(T_v) - \mathit{misclass}(v).
\]
If $\mathit{Gain\_prune}(v) > 0$, pruning reduces the validation error. If it is zero, the simpler tree is often preferred. Starting from the fully grown tree, we evaluate the possible pruning operations bottom-up, prune the node with the largest positive gain and repeat until all the remaining gains are negative.

\subsubsection{Interpretability}
Generally pruning improves the tree but does not solve everything. Some structural problems may remain:
\begin{itemize}
  \item \g{repetition}: the same attribute is tested multiple times along a single root-to-leaf path, which makes the path longer and less readable;
  \item \g{replication}: similar or identical subtrees appear in different parts of the tree, which increases its size without adding much explanatory value.
\end{itemize}
As alternatives, we can use multivariate splits, or extract IF-THEN rules from the tree and prune or simplify the rules directly.

Decision trees are considered very interpretable classifiers, because we can transform the tree into rules: each root-to-leaf path becomes an IF-THEN rule. We can have \g{global interpretability}, related to the overall structure of the tree (we want it not too big), and \g{local interpretability}, where for a specific instance we can follow its unique path in the tree to understand the prediction.

\subsection{Bayesian classifier}

A Bayesian classifier is a statistical classifier: instead of returning only a label, it returns for every class the probability that the instance belongs to that class. As classification criterion we choose the class with the highest probability. One problem is that we have to compute many probabilities, and this could be heavy. We will see the \g{naive Bayes classifier}.

Notation (slide 89):
\begin{center}
\begin{tabular}{@{}ll@{}}
  \toprule
  Symbol & Meaning \\
  \midrule
  $X$ & observed data tuple (the evidence), class label unknown \\
  $H$ & hypothesis: $X$ belongs to a class $C$ \\
  $P(H)$ & prior probability of the hypothesis \\
  $P(X)$ & evidence: probability that $X$ is observed \\
  $P(X \mid H)$ & likelihood: probability of observing $X$ if $H$ holds \\
  $P(H \mid X)$ & posterior: probability that $H$ holds after observing $X$ \\
  \bottomrule
\end{tabular}
\end{center}
By Bayes' theorem $P(C_i \mid X) = P(X \mid C_i)\,P(C_i)/P(X)$, and we predict the class $C_i$ with the largest posterior. Since $P(X)$ is the same for all the classes, it is enough to maximize $P(X \mid C_i)\,P(C_i)$.

To reduce the number of probabilities to compute we use the \g{naive assumption}: given the class $C_i$, the attributes are conditionally independent, so
\[
  P(X \mid C_i) = \prod_{k=1}^{n} P(x_k \mid C_i).
\]
The assumption is useful even when it is not perfectly true: it turns a hard joint-probability estimation problem into several reliable class-conditioned frequency estimates.

The naive Bayes classifier can have a problem: a single zero conditional probability makes the whole product zero. To solve this we use the \g{Laplacian correction}: we add one pseudo-count to each possible value, so that no probability is zero.

Advantages:
\begin{itemize}
  \item easy to implement;
  \item often gives good results;
  \item computationally efficient, because it mainly requires class-conditional counts.
\end{itemize}
Limitations:
\begin{itemize}
  \item it assumes class-conditional independence, while real variables are often dependent.
\end{itemize}

\subsection{Bayesian belief networks}

When the independence assumption does not hold we need richer models, such as \g{Bayesian belief networks}. They are defined by a directed acyclic graph and by \g{conditional probability tables} (CPT): nodes are random variables and links represent dependencies. The graph encodes conditional independencies that simplify the joint distribution: each variable is conditionally independent of its non-descendants given its parents, so
\[
  P(X_1,\dots,X_n) = \prod_{i=1}^{n} P(X_i \mid \mathit{Parents}(X_i)).
\]

For classification, any node can be selected as the output node, representing the class label (there may also be more than one output node), and the classifier can return a probability distribution over the classes.
